\section{Introduction}\label{sec:introduction}

Coverage is a basic operation in inspection, infrastructure monitoring, environmental sensing, and search without reliable target priors. A coverage system must distribute finite sensing resources over a spatial domain, decide the order in which regions are visited, and determine when a visit can be counted. Classical coverage planning represents the domain by cells, waypoints, or continuous density functions and then optimizes a geometric surrogate such as travel distance, makespan, or accumulated sensing cost \cite{cortes2004coverage,galceran2013survey,tan2021review}. Multi-robot extensions exploit parallel motion and task allocation to reduce completion time. Recent methods further incorporate local-search improvement, online assignment, concurrent planning and execution, and priority-dependent first-visit delay \cite{wang2024apfcpp,wang2025macplanner,mitra2024ondemand,mitra2025concurrent,lee2026priority}. These developments make the planned routes increasingly responsive, but the state variable used by most planners remains an abstract visit: a node is treated as covered when the route reaches it or when the scheduled action is assumed to finish.

The abstraction is less reliable for networked directional cameras. Such a camera covers space by changing its optical axis rather than translating its base. A single pointing direction may cover many cells, and its value depends on field-of-view geometry, camera placement, and overlap with other viewpoints. The time needed to establish an observation is also more than angular displacement divided by a nominal speed. It includes command delivery, device-dependent and direction-dependent gimbal motion, settling, the prescribed dwell, video refresh, and the association of a usable image with the corresponding pose. A controller may acknowledge a velocity command while the camera is still moving. A pose may satisfy its angular tolerance while the available stream still contains a frame acquired before arrival. Conversely, a fresh frame cannot support spatial coverage if its pose association is absent or too old. The physical observation event therefore lies downstream of the planned viewpoint and downstream of the control command.

This distinction matters for both optimization and evaluation. If coverage is updated at the scheduled arrival time, the objective is calculated from events that may not yet have happened. The resulting first-coverage curve can be optimistic even when every waypoint is eventually reached. If a device becomes temporarily unavailable, immediately replacing its queue with a newly optimized queue introduces a second problem: the replacement may be worse than simply allowing healthy devices to continue their current work. A repair mechanism is useful only when its predicted benefit exceeds the cost and uncertainty of switching, and its claimed effect must be visible in subsequent physical observations. These are not isolated implementation details. Together they form a planning--execution mismatch: the planner reasons over a set of feasible viewpoint--cell relations and predicted completion times, whereas the hardware exposes commands, delayed pose samples, image frames, and availability events.

Active camera networks have addressed high-speed tracking, target handoff, predictive scheduling, and zoom control \cite{e2024wisecam,jiang2021pan,zhang2026capture}. Robust active-perception and resilient allocation studies have separately examined motion uncertainty, perception quality, failures, and temporal constraints \cite{xanthidis2024resivis,choudhury2022dynamic,notomista2022resilient}. Yet these lines of work do not by themselves provide a coverage state that is traceable from a planned pointing command to a measured pose and a newly acquired frame, nor do they require every accepted repair to dominate an executable continuation baseline under the same coverage objective. The missing link is an operational contract between planning, physical execution, and observation evidence.

This article formulates that contract for first effective coverage of a fixed three-dimensional region by networked directional cameras. The region is discretized into equal-weight evaluation cells. The primary metric is the fixed-horizon area under the fraction of uncovered cells, termed the uncovered-area criterion (UAC). UAC accounts for the first observation time of every cell rather than only the last completion event. Complete-coverage time, T100, is imposed as a tail constraint instead of being combined with UAC through a problem-dependent weighted sum. Under the limited interpretation of a stationary target whose initial cell is uniformly distributed and whose first valid observation implies detection, UAC equals a truncated expected detection delay. Without those assumptions it is, more precisely, a truncated first-observation latency and is reported as such throughout the article.

The resulting method, ECR-Cover, links three components. First, initial planning retains a complete motion-aware greedy (MAG) route as a fallback and uses calibrated end-to-end execution-time upper bounds to evaluate deterministic local modifications. A candidate can replace the incumbent only if it preserves complete feasible coverage, strictly reduces robust UAC, and respects the MAG T100 bound. Second, execution uses a state machine in which an acknowledged command, a continuously stable measured pose, a dwell interval, a usable fresh frame, and bounded pose--frame time separation are all required before coverage is recorded. The recorded set is generated by projecting the measured pose associated with that frame; the planned pose never writes coverage directly. Third, temporary unavailability triggers residual-task repair. The already confirmed set is frozen, active tasks on healthy devices are retained, and the candidate repair is compared with the executable policy of continuing the original queues. Repair is committed only when it improves predicted residual UAC beyond a switching threshold, does not increase predicted residual T100, and contains a predicted reassigned observation before the unavailable device can recover.

The contribution is consequently a closed-loop coverage method rather than a new general-purpose optimization theorem. Its formal results are structural properties of the frozen planning model and the implemented evidence rules: (i) the initial planner preserves the MAG baseline in robust UAC and T100; (ii) every logged coverage update has a matching command--pose--frame evidence chain; (iii) coverage completes under explicit reachability, recovery, and eventual-confirmation conditions; and (iv) an accepted repair preserves its original-queue baseline in the predicted residual problem. Hardware results then test whether these model-level properties have useful observable consequences. Four ZR10 pan--tilt cameras supplied 240 stratified transitions for execution-time modeling and 70 runs under frozen nominal, interruption, and repair-ablation protocols. The experiments evaluate first-coverage latency, completion tails, physical evidence consistency, and computation time without using a target detector or target ground truth.

